\documentclass[12pt]{article}
%^ Start of document
\usepackage{times}  %Makes text TimesNewRoman
\usepackage{setspace} %Allows you to set single or double space 
\usepackage{biblatex} %Imports biblatex package
\usepackage{graphicx} % Required for inserting images
\usepackage{authblk}
\usepackage[colorlinks=true, urlcolor=blue, linkcolor=blue]{hyperref}
\usepackage{subfig}
\usepackage{float}
\usepackage{tabularx}
\usepackage{multirow}
\usepackage[paperheight=11in,
   paperwidth=8.5in,
   top=20mm,
   bottom=20mm,
   left=40mm,
   right=40mm]{geometry}
\usepackage{moresize}
\usepackage{tikz}
\usepackage{xcolor}
\usepackage{physics}
\usepackage{amssymb}
\usepackage{mathrsfs}
\usepackage{amsmath}
\usepackage{ragged2e}
\usepackage{comment}
\usepackage{xcolor}
\addbibresource{seminar.bib}
\begin{document}
\singlespacing  %Sets single spacing for whole document
\title{Conducting Polymers - Edman}

\author[1]{Zachary Tullier}
\affil[1]{Student\\Physics Department\\ University of Houston - Clear Lake \\Houston, Texas\\U.S}
\date{}
\maketitle

\section{Introduction to the Speaker and Background}
    \textbf{Speaker:} Dr. Jennifer Edman, Dean of the College of Science, is an experienced researcher in materials science with expertise in conducting polymers. Her career spans academia and industry, with significant contributions to organic chemistry and material applications. Dr. Edman has been involved in research on conducting polymers since the early 1990s. She transitioned her focus from fundamental studies to practical applications such as sensors, water purification, and biomedical devices.


\section{Introduction to Conducting Polymers}
    Conducting polymers are electro-active materials with intrinsic electrical conductivity due to their unique electronic structures based on alternating single and double bonds. Reversibility is a feature of polimers where they can switch between insulating and conducting states for millions of cycles.
    The chemical properties of the polymers can be tailored for specific applications. Each polymer type has a wide range of uses in electronics, energy, and healthcare. The discovery of conducting polymers was awarded the 2000 Nobel Prize in Chemistry to Heeger, MacDiarmid, and Shirakawa, highlighting its significance in materials science.

\section{Fundamental Properties of Conducting Polymers}
    Alternating single and double bonds form delocalized pi-electron clouds, enabling unique electronic and optical behaviors. The conjugation length influences conductivity. The energy gap between the highest occupied molecular orbital (HOMO) and the lowest unoccupied molecular orbital (LUMO) determines conductivity. Smaller band gaps improve conductivity and broaden application possibilities. Conducting polymers are often rigid and insoluble. Chemical modifications, such as adding flexible side chains, improve solubility and enable practical applications.

\section{Applications of Conducting Polymers}
    Conducting Polymers can be used in photovoltaics (Organic solar cells for efficient sunlight harvesting), pseudo-capacitors for energy storage, or protective coatings for static dissipation. They can be used as biosensors, detecting genetic markers for diseases or enabling scaffolds for biomedical applications. In the industrial worlds they can be used in smart windows, glasses, and adaptive surfaces or offer rapid tinting control through electric fields. Photocatalytic membranes made from conducting polymers are being developed for pollutant degradation, providing sustainable clean water solutions.

\section{Electroactive Polymers and Practical Use Cases}
    Conducting polymers exhibit properties that change in response to electric fields, enabling dynamic stability and adaptability. Properly designed polymers endure millions of cycles, making them reliable for long-term applications. Photocatalytic water purification is a key area of ongoing research. Conducting polymers can serve as membranes to degrade pollutants and enhance water treatment processes.


\section{Structural and Chemical Modifications}
    Incorporating heteroatoms like sulfur, nitrogen, or oxygen improves oxidation potential and stability. Adding electron-donating groups lowers the band gap. Flexible side chains enhance solubility and processability.

\section{Advances in Fabrication: From Films to Nanofibers}
    Thin films are widely used in sensors, energy storage devices, and photovoltaic cells. Nanofibers produced via electrospinning provide high surface areas. With alternative fabrication processes, researches often have to trade off stability vs. switching speed and solubility vs. performance. Counterions and electrolytes impact conductivity and stability. Even in small scale experiments ensuring quality and consistency is difficult, but for this to be put in production, the quality must be ensured for larger batches. This comes with its own unique challenges. 


\section{Personal and Historical Insights}
Dr. Edman shared anecdotes from her career, including interactions with Nobel laureates and the evolution of conducting polymer research. She highlighted the shift from fundamental discoveries to real-world applications. Dr. Edman emphasized the importance of collaboration across chemistry, physics, and materials science to address global challenges.

\section{Conclusion}
Conducting polymers represent a trans-formative material in modern science. From energy storage to biomedical applications, their versatility and adaptability drive innovation. Dr. Edman’s work highlights the importance of interdisciplinary research in unlocking the full potential of these materials.

\end{document}

